% Sección 8
\section{Lo que no sabemos todavía y a quién preguntarlo}

Las preguntas abiertas se responden mejor con una persona que estuvo en la operación que con más prensa. Una respuesta de alguien de gestión del riesgo o de un organismo de socorro, citada con nombre y cargo, pesa más que cualquier cifra de este informe. La tabla~\ref{tab:vacios} pone primero los vacíos que más afectan al despacho. La columna de la derecha propone a quién preguntar; cuando la investigación no lo dice, la propuesta es nuestra.

{\footnotesize
\setlength{\tabcolsep}{4pt}
\renewcommand{\arraystretch}{1.3}
\begin{longtable}{L{0.025\linewidth} L{0.35\linewidth} L{0.28\linewidth} L{0.26\linewidth}}
\caption{Vacíos de información y a quién preguntar.}\label{tab:vacios}\\
\toprule
\textbf{\#} & \textbf{Lo que no sabemos} & \textbf{Por qué importa} & \textbf{A quién preguntarlo}\\
\midrule
\endfirsthead
\multicolumn{4}{l}{\small\itshape Tabla~\ref{tab:vacios}, continuación}\\
\toprule
\textbf{\#} & \textbf{Lo que no sabemos} & \textbf{Por qué importa} & \textbf{A quién preguntarlo}\\
\midrule
\endhead
\midrule
\multicolumn{4}{r}{\small\itshape Sigue en la página siguiente}\\
\endfoot
\bottomrule
\endlastfoot
1 & Cómo despacha hoy, paso a paso, el Banco de Alimentos de Manizales, quién lo dirige y qué inventario y comités atiende & Es el usuario principal que propone el equipo, y nadie ha hablado con él & Banco de Alimentos de Manizales; ABACO\\
2 & Qué vías de acceso se cerraron por el sismo, desde cuándo y hasta cuándo, incluidas La Siria y las salidas a Medellín y a Bogotá & Solo hay prensa de los dos primeros días, y sin la duración no sabemos cuánto pesaron los cierres & UGR de Manizales; Policía de Tránsito; INVÍAS; Gobernación de Caldas\\
3 & Si hubo desabastecimiento de alimentos en las primeras 48 horas, y en qué barrios & La ficha del reto lo afirma y no encontramos fuente & Alguien de gestión del riesgo presente en la hackathon; UGR de Manizales\\
4 & Si se dañaron o vencieron alimentos donados en Manizales, y cuántos & Es el segundo indicador del reto & ABACO; Banco de Alimentos de Manizales; Cruz Roja de Caldas\\
5 & Cuántas raciones, kits y familias se atendieron en Manizales, y desde qué centros de acopio & Solo hay cifras nacionales de ABACO y de donantes privados & Gestora social de Manizales (acopio del Coliseo Menor); ABACO\\
6 & Qué comunas y barrios concentraron a los damnificados y cuáles quedaron aislados & Define dónde están los puntos críticos & UGR de Manizales\\
7 & Cuántas familias hubo en albergues y cómo cambió esa cifra semana a semana & Hay cifras de personas y no de familias & UGR de Manizales; Alcaldía\\
8 & Cómo operan los comités barriales y los transportadores, y qué datos tienen & Son quienes piden y quienes llevan, y no hay información pública & Juntas de acción comunal; transportadores; Alcaldía (consolidado de juntas)\\
9 & Si hubo embotellamientos críticos en los accesos a la ciudad & La ficha los menciona y solo hay antecedentes de 2021 a 2024 & Policía de Tránsito; UGR de Manizales\\
10 & Qué dicen la Cruz Roja de Caldas, la Defensa Civil de Caldas, la UDEGER y la UGR de Manizales, y quiénes las dirigen & No encontramos declaraciones propias & Cada entidad, de forma directa\\
11 & Cuántos pedidos llegan mal ubicados & La simulación supone un 10~\% sin fuente & Comités barriales; Banco de Alimentos de Manizales\\
12 & Cuántos pedidos tiene una jornada real & La ventaja simulada va de 12 a 46 puntos según la carga & Banco de Alimentos de Manizales (datos de despachos de agosto)\\
13 & Si existen ERIS-UNAL, el software de ruteo adaptativo y los datos de la ficha 04 & El banco de capacidades del sitio no tiene rastro público de ellos & Mesa de la hackathon\\
14 & Si la formulación del ruteo y el papel del despachador son razonables & Da respaldo académico local & William Sarache, Diana Cárdenas o Juan Camilo López Vargas, investigadores de la Sede Manizales en logística humanitaria\\
15 & Cómo cambiaron los precios de los alimentos en Manizales & Solo hay datos nacionales y el boletín de Centro Galerías & DANE (SIPSA); Fenalco Caldas\\
16 & La población oficial de Manizales & La cifra viene de una página que reproduce al DANE & DANE\\
17 & Qué alternativas usan hoy las entidades y qué marco oficial rige la recuperación & Responde a dos preguntas de la mentoría & \eninvestigacion\\
18 & Qué dijeron los actores en redes, prensa e informes de gestión & Sustituye en parte a las entrevistas que faltan & \eninvestigacion\\
\end{longtable}
}
